% GENERATED FILE. Edit canonical lessons, then run npm run generate:books.
% canonical-source-hash: fnv1a64:553ae5336bedd019
% canonical-lessons: SA-C478-vedana, SA-C478-kasate, SA-C478-atisarah, SA-C478-klesah, SA-C478-asahyah

\chapter{Aches and Pains}
\label{ch:sa-aches-and-pains}
\hlchaptermodality{478}

\begin{chapteropening}
\textbf{By the end of this chapter:} \emph{I can say where it hurts and how badly.}

Complete the last lesson of chapter 478: I can say where it hurts and how badly.
\end{chapteropening}

\section[\texorpdfstring{\sk{वेदना} (vedanā)}{वेदना (vedanā)}]{\sk{वेदना} (vedanā) — agony; acute pain}

\label{lesson:SA-C478-vedana}



\begin{quote}
Before the new one: say the Sanskrit for a burning feeling, then the Sanskrit for runs a fever.
\end{quote}

\subsection*{You'll want to know: \sk{वेदना}}
\textbf{\sk{वेदना}} — \emph{vedanā} — \textquotedblleft{}agony; acute pain\textquotedblright{}.

\begin{quote}
\sk{रात्रौ} \sk{मम} \sk{दन्ते} \sk{वेदना} \sk{अभवत्।}
\end{quote}

\begin{quote}
\emph{rātrau mama dante vedanā abhavat.}
\end{quote}

\begin{quote}
In the night I had a terrible toothache.
\end{quote}

A feminine noun.

\begin{cousinweb}
\textbf{\sk{वेदना}} (\emph{vedanā}) comes from the root √vid, \textquotedblleft{}to know, to feel\textquotedblright{}: what is felt; MW gives \textquotedblleft{}pain, torture, agony\textquotedblright{} (MW).
\end{cousinweb}

\begin{grammarlens}[title={What you've built}]
One more word for describing what you lived and felt.
\end{grammarlens}

\subsection*{Your turn}
\begin{itemize}
\raggedright
  \item \emph{Say it:} \emph{vedanā}
  \item \emph{Say it:} \emph{vedanā}, once more
  \item \emph{Say it:} \emph{rātrau mama dante vedanā abhavat.} — In the night I had a terrible toothache.
  \item \emph{Recall:} say the Sanskrit for a burning feeling, then the Sanskrit for runs a fever, then say \emph{vedanā} again
\end{itemize}

\begin{tcolorbox}[breakable,colback=teal!4,colframe=teal!35!black,title={Before you move on}]
What does \sk{वेदना} mean? (\textquotedblleft{}Agony; acute pain\textquotedblright{}.) Say it once more.
\end{tcolorbox}
\section[\texorpdfstring{\sk{कासते} (kāsate)}{कासते (kāsate)}]{\sk{कासते} (kāsate) — coughs}

\label{lesson:SA-C478-kasate}



\begin{quote}
Before the new one: say the Sanskrit for runs a fever, then the Sanskrit for agony.
\end{quote}

\subsection*{You'll want to know: \sk{कासते}}
\textbf{\sk{कासते}} — \emph{kāsate} — \textquotedblleft{}coughs\textquotedblright{}.

\begin{quote}
\sk{शीतेन} \sk{बालः} \sk{रात्रौ} \sk{कासते।}
\end{quote}

\begin{quote}
\emph{śītena bālaḥ rātrau kāsate.}
\end{quote}

\begin{quote}
In the cold the child coughs at night.
\end{quote}

Past (he …): \sk{सः} \sk{रात्रौ} \sk{कासितवान्}, \emph{saḥ rātrau kāsitavān}.

\begin{cousinweb}
\textbf{\sk{कासते}} (\emph{kāsate}) comes from the root √kās, \textquotedblleft{}to cough\textquotedblright{}, the root of \textbf{\sk{कासः}} (MW).
\end{cousinweb}

\begin{grammarlens}[title={What you've built}]
One more word for describing what you lived and felt.
\end{grammarlens}

\subsection*{Your turn}
\begin{itemize}
\raggedright
  \item \emph{Say it:} \emph{kāsate}
  \item \emph{Say it:} \emph{kāsate}, once more
  \item \emph{Say it:} \emph{śītena bālaḥ rātrau kāsate.} — In the cold the child coughs at night.
  \item \emph{Recall:} say the Sanskrit for runs a fever, then the Sanskrit for agony, then say \emph{kāsate} again
\end{itemize}

\begin{tcolorbox}[breakable,colback=teal!4,colframe=teal!35!black,title={Before you move on}]
What does \sk{कासते} mean? (\textquotedblleft{}Coughs\textquotedblright{}.) Say it once more.
\end{tcolorbox}
\section[\texorpdfstring{\sk{अतिसारः} (atisāraḥ)}{अतिसारः (atisāraḥ)}]{\sk{अतिसारः} (atisāraḥ) — dysentery}

\label{lesson:SA-C478-atisarah}



\begin{quote}
Before the new one: say the Sanskrit for agony, then the Sanskrit for coughs.
\end{quote}

\subsection*{You'll want to know: \sk{अतिसारः}}
\textbf{\sk{अतिसारः}} — \emph{atisāraḥ} — \textquotedblleft{}dysentery\textquotedblright{}.

\begin{quote}
\sk{मलिनेन} \sk{जलेन} \sk{अतिसारः} \sk{भवति।}
\end{quote}

\begin{quote}
\emph{malinena jalena atisāraḥ bhavati.}
\end{quote}

\begin{quote}
Dirty water causes dysentery.
\end{quote}

A masculine noun.

\begin{cousinweb}
\textbf{\sk{अतिसारः}} (\emph{atisāraḥ}) is built from \emph{ati}, \textquotedblleft{}beyond\textquotedblright{}, and the root √sṛ, \textquotedblleft{}to flow\textquotedblright{}: a flowing over; MW gives \textquotedblleft{}purging, dysentery\textquotedblright{} (MW).
\end{cousinweb}

\begin{grammarlens}[title={What you've built}]
One more word for describing what you lived and felt.
\end{grammarlens}

\subsection*{Your turn}
\begin{itemize}
\raggedright
  \item \emph{Say it:} \emph{atisāraḥ}
  \item \emph{Say it:} \emph{atisāraḥ}, once more
  \item \emph{Say it:} \emph{malinena jalena atisāraḥ bhavati.} — Dirty water causes dysentery.
  \item \emph{Recall:} say the Sanskrit for agony, then the Sanskrit for coughs, then say \emph{atisāraḥ} again
\end{itemize}

\begin{tcolorbox}[breakable,colback=teal!4,colframe=teal!35!black,title={Before you move on}]
What does \sk{अतिसारः} mean? (\textquotedblleft{}Dysentery\textquotedblright{}.) Say it once more.
\end{tcolorbox}
\section[\texorpdfstring{\sk{क्लेशः} (kleśaḥ)}{क्लेशः (kleśaḥ)}]{\sk{क्लेशः} (kleśaḥ) — suffering}

\label{lesson:SA-C478-klesah}



\begin{quote}
Before the new one: say the Sanskrit for coughs, then the Sanskrit for dysentery.
\end{quote}

\subsection*{You'll want to know: \sk{क्लेशः}}
\textbf{\sk{क्लेशः}} — \emph{kleśaḥ} — \textquotedblleft{}suffering\textquotedblright{}.

\begin{quote}
\sk{रोगः} \sk{बहु} \sk{क्लेशं} \sk{ददाति।}
\end{quote}

\begin{quote}
\emph{rogaḥ bahu kleśaṃ dadāti.}
\end{quote}

\begin{quote}
Illness causes much suffering.
\end{quote}

A masculine noun.

\begin{cousinweb}
\textbf{\sk{क्लेशः}} (\emph{kleśaḥ}) comes from the root √kliś, \textquotedblleft{}to torment, to suffer\textquotedblright{}; MW gives \textquotedblleft{}pain, affliction, distress, pain from disease\textquotedblright{} (MW).
\end{cousinweb}

\begin{grammarlens}[title={What you've built}]
One more word for describing what you lived and felt.
\end{grammarlens}

\subsection*{Your turn}
\begin{itemize}
\raggedright
  \item \emph{Say it:} \emph{kleśaḥ}
  \item \emph{Say it:} \emph{kleśaḥ}, once more
  \item \emph{Say it:} \emph{rogaḥ bahu kleśaṃ dadāti.} — Illness causes much suffering.
  \item \emph{Recall:} say the Sanskrit for coughs, then the Sanskrit for dysentery, then say \emph{kleśaḥ} again
\end{itemize}

\begin{tcolorbox}[breakable,colback=teal!4,colframe=teal!35!black,title={Before you move on}]
What does \sk{क्लेशः} mean? (\textquotedblleft{}Suffering\textquotedblright{}.) Say it once more.
\end{tcolorbox}
\section[\texorpdfstring{\sk{असह्यः} (asahyaḥ)}{असह्यः (asahyaḥ)}]{\sk{असह्यः} (asahyaḥ) — unbearable}

\label{lesson:SA-C478-asahyah}



\begin{quote}
Before the new one: say the Sanskrit for dysentery, then the Sanskrit for suffering.
\end{quote}

\subsection*{You'll want to know: \sk{असह्यः}}
\textbf{\sk{असह्यः}} — \emph{asahyaḥ} — \textquotedblleft{}unbearable\textquotedblright{}.

\begin{quote}
\sk{तस्य} \sk{वेदना} \sk{असह्या} \sk{अस्ति।}
\end{quote}

\begin{quote}
\emph{tasya vedanā asahyā asti.}
\end{quote}

\begin{quote}
His pain is unbearable.
\end{quote}

\begin{cousinweb}
\textbf{\sk{असह्यः}} (\emph{asahyaḥ}) is built from \emph{a}, \textquotedblleft{}not\textquotedblright{}, and \emph{sahya}, \textquotedblleft{}to be borne\textquotedblright{}, from the root √sah of \textbf{\sk{सहते}}; MW gives \textquotedblleft{}unbearable, insufferable\textquotedblright{} (MW).
\end{cousinweb}

\begin{grammarlens}[title={What you've built}]
One more word for describing what you lived and felt.
\end{grammarlens}

\subsection*{Your turn}
\begin{itemize}
\raggedright
  \item \emph{Say it:} \emph{asahyaḥ}
  \item \emph{Say it:} \emph{asahyaḥ}, once more
  \item \emph{Say it:} \emph{tasya vedanā asahyā asti.} — His pain is unbearable.
  \item \emph{Recall:} say the Sanskrit for dysentery, then the Sanskrit for suffering, then say \emph{asahyaḥ} again
\end{itemize}

\begin{tcolorbox}[breakable,colback=teal!4,colframe=teal!35!black,title={Before you move on}]
What does \sk{असह्यः} mean? (\textquotedblleft{}Unbearable\textquotedblright{}.) Say it once more.
\end{tcolorbox}
